\chapter{Hipótesis y Objetivos del Proyecto}

\section{Hipótesis}

El diseño y la integración de una estación inteligente de control ESD de bajo
costo permitirán transformar la verificación manual de pulseras antiestáticas en
un evento trazable y auditable. Para ello, el sistema integrará un probador
comercial HZR-171, mecanismos de identificación de usuarios, una interfaz
gráfica táctil y almacenamiento local de los registros generados durante cada
prueba.

El prototipo generará registros completos en al menos el 95\% de los ciclos de
prueba ejecutados. Cada registro deberá incluir el nombre y el identificador del
usuario, la fecha, la hora, el resultado ESD, la temperatura y la humedad relativa.
Además, el costo estimado de materiales se mantendrá por debajo de USD~350.

\section{Objetivos}

\subsection{Objetivo general}

Diseñar, implementar y validar un prototipo de estación inteligente de bajo
costo para el control y la trazabilidad de verificaciones ESD en laboratorios de
electrónica, mediante la integración de un probador de pulseras antiestáticas,
mecanismos de identificación de usuarios, una interfaz gráfica y almacenamiento
local de registros.

\subsection{Objetivos específicos}

\begin{itemize}

\item Diseñar la arquitectura electrónica, de comunicación, alimentación e 
integración física de la estación, con base en los requerimientos 
funcionales y no funcionales del sistema.

\item Diseñar e implementar los subsistemas de identificación de usuarios y
adquisición del resultado de la prueba ESD, mediante la integración del lector
de código de barras, el reconocimiento facial complementario y las señales del
probador HZR-171.

\item Diseñar e implementar la interfaz gráfica que guíe al usuario durante el
ciclo de prueba y el sistema de almacenamiento de los resultados y generación
de reportes históricos trazables.

\item Validar el desempeño funcional del prototipo mediante pruebas individuales
de los subsistemas y pruebas de integración del sistema completo, estas últimas
realizadas en un laboratorio de electrónica, utilizando los indicadores de
trazabilidad e integración funcional definidos para el proyecto.

\end{itemize}

\section{Actividades y entregables por objetivo}

\subsection{Objetivo específico 1}

\textbf{Diseñar la arquitectura electrónica, de comunicación, alimentación e 
integración física de la estación, con base en los 
requerimientos funcionales y no funcionales del sistema.}

\paragraph{Actividades}

\begin{itemize}
\item Analizar los requerimientos funcionales y 
no funcionales del sistema como base para el diseño del prototipo.
\item Seleccionar y justificar los componentes electrónicos de acuerdo con los
requerimientos del proyecto.
\item Definir la arquitectura general y las interfaces de comunicación entre
los subsistemas.
\item Diseñar el sistema de alimentación eléctrica y la distribución de energía
del prototipo.
\item Definir la asignación de pines y seleccionar los 
conectores necesarios para la interconexión de los subsistemas.
\item Elaborar los diagramas eléctricos y de interconexión del sistema.
\item Diseñar las PCB necesarias para integrar los módulos electrónicos.
\item Fabricar y verificar eléctricamente las PCB desarrolladas.
\item Diseñar el modelo CAD y la distribución física de los componentes.
\item Fabricar la estructura y ensamblar el prototipo.
\end{itemize}

\paragraph{Entregables}

\begin{itemize}
\item Especificación de requerimientos y criterios de diseño.
\item Arquitectura electrónica y diagramas de bloques e interconexión.
\item Diseño del sistema de alimentación y distribución de pines.
\item PCB diseñadas, fabricadas y verificadas.
\item Modelo CAD y prototipo físico ensamblado.
\end{itemize}


\subsection{Objetivo específico 2}

\textbf{Diseñar e implementar los subsistemas de identificación de usuarios y
adquisición del resultado de la prueba ESD, mediante la integración del lector
de código de barras, el reconocimiento facial complementario y las señales del
probador HZR-171.}

\paragraph{Actividades}

\begin{itemize}
\item Definir el flujo de identificación y su relación con el ciclo de prueba.
\item Diseñar e implementar la comunicación con el lector de código de barras
GM67.
\item Desarrollar las rutinas de lectura y validación de identificadores.
\item Diseñar e implementar el proceso complementario de registro y
reconocimiento facial.
\item Diseñar la adquisición de las señales \textit{LO}, \textit{OK} y
\textit{HI} del probador HZR-171.
\item Integrar el probador HZR-171 con el controlador secundario y el
controlador principal.
\item Asociar la identificación del usuario con el resultado de la prueba ESD.
\item Verificar funcionalmente los subsistemas de identificación y adquisición.
\end{itemize}

\paragraph{Entregables}

\begin{itemize}
\item Diseño documentado del flujo de identificación y adquisición ESD.
\item Subsistema de identificación por código de barras funcional.
\item Subsistema complementario de reconocimiento facial.
\item Adquisición funcional de los estados del probador HZR-171.
\item Resultados de las pruebas funcionales de los subsistemas de identificación
y adquisición del resultado ESD.
\end{itemize}

\subsection{Objetivo específico 3}

\textbf{Diseñar e implementar la interfaz gráfica que guíe al usuario durante
el ciclo de prueba y el sistema de almacenamiento de los resultados y
generación de reportes históricos trazables.}

\paragraph{Actividades}

\begin{itemize}

\item Definir el flujo de interacción y diseñar la estructura de la interfaz
gráfica.

\item Implementar la interfaz gráfica y su flujo de navegación para guiar al
usuario durante el ciclo de prueba y presentar su información de identificación
y el resultado ESD.

\item Diseñar la estructura de almacenamiento de los registros.

\item Integrar el sensor SHT31-D como fuente de información complementaria de
temperatura y humedad relativa para cada registro.

\item Implementar el almacenamiento local del nombre y el identificador del
usuario, la fecha, la hora, el resultado ESD y las variables ambientales asociadas.

\item Implementar la consulta de los registros almacenados y la generación de
reportes históricos a partir de estos.

\item Verificar el funcionamiento de la interfaz, el almacenamiento, la consulta
de registros y la generación de reportes.

\end{itemize}

\paragraph{Entregables}

\begin{itemize}

\item Interfaz gráfica táctil funcional.

\item Diseño documentado del flujo de interacción y de la estructura de datos.

\item Sistema local de almacenamiento, consulta de registros y generación de
reportes históricos.

\item Resultados de las pruebas funcionales de interacción y gestión de
información.

\end{itemize}

\subsection{Objetivo específico 4}

\textbf{Validar el desempeño funcional del prototipo mediante la evaluación de
los resultados de las pruebas individuales de los subsistemas y la ejecución de
pruebas de integración del sistema completo en un laboratorio de electrónica,
utilizando los indicadores de trazabilidad e integración funcional definidos
para el proyecto.}

\paragraph{Actividades}

\begin{itemize}
\item Definir el procedimiento de validación y los criterios de aceptación del
prototipo.
\item Evaluar los resultados de las pruebas funcionales individuales de los
subsistemas desarrollados en los objetivos anteriores.
\item Ejecutar en el laboratorio TEM de Intel Costa Rica las pruebas de
integración del sistema completo, desde la identificación del usuario hasta el
almacenamiento del registro, incluida la consulta de la información y la
generación de reportes históricos.
\item Si circunstancias fuera del control del proyecto impiden realizar las
pruebas de integración en TEM, aplicar el mismo procedimiento en un laboratorio
de la Escuela de Ingeniería Electrónica del Tecnológico de Costa Rica, Campus
Tecnológico Local San Carlos.
\item Calcular el índice de trazabilidad \(TZ\) y la tasa de integración
funcional \(IF\) a partir de los resultados de las pruebas de integración.
\item Analizar los resultados obtenidos y determinar el cumplimiento de los
criterios de aceptación establecidos.
\item Documentar las incidencias, limitaciones y conclusiones de la validación
del prototipo.
\end{itemize}

\paragraph{Entregables}

\begin{itemize}
\item Procedimiento de validación y criterios de aceptación.
\item Evaluación consolidada de los resultados de las pruebas funcionales
individuales de los subsistemas.
\item Resultados documentados de las pruebas de integración del sistema
completo.
\item Valores calculados para el índice de trazabilidad \(TZ\) y la tasa de
integración funcional \(IF\).
\item Informe de validación del prototipo, con el análisis de cumplimiento, las
incidencias, las limitaciones y las conclusiones obtenidas.
\end{itemize}

\section{Alcances y limitaciones del proyecto}

\subsection{Alcances}

El resultado del proyecto será un prototipo académico funcional de una estación
de control ESD para pulseras antiestáticas, capaz de asociar cada prueba con la
identificación del usuario y conservar localmente la información necesaria para
su consulta. El prototipo se entregará ensamblado en una estructura física y se
validará mediante pruebas de los subsistemas y del funcionamiento integrado en
un entorno de laboratorio. Su alcance corresponde a una demostración funcional
de la solución propuesta y no a un producto comercial certificado. Tampoco se
contempla su integración con bases de datos, sistemas de control de acceso u
otras plataformas empresariales. Las funciones y condiciones esperadas del
prototipo se detallan en los requisitos funcionales y no funcionales de la
sección~\ref{sec:requisitos_sistema}.

\subsection{Limitaciones}

El proyecto estará sujeto a las siguientes limitaciones:

\begin{itemize}
\item La verificación ESD dependerá del rango de operación y de los criterios de
aprobación del probador comercial HZR-171. El prototipo se limitará a pulseras
antiestáticas y no evaluará calzado ni tobilleras ESD.
\item El reconocimiento facial podrá producir falsos positivos o falsos
negativos y su desempeño dependerá, entre otros factores, de la iluminación, la
posición del usuario y la calidad de la imagen. Las evaluaciones de NIST
muestran que las tasas de error varían entre algoritmos y conjuntos de imágenes
\cite{NISTFRTE11}. Por ello, no se establece de antemano una tasa fija de
acierto, como 95~\%. Esta función se evaluará experimentalmente y se mantendrá
como método complementario; la lectura del código de barras será el mecanismo
principal de identificación.
\item El almacenamiento será local y no se contempla integración con servicios
en la nube, bases de datos corporativas, sistemas de control de acceso ni
plataformas de gestión de personal.
\item Las pruebas de integración están planificadas en el laboratorio TEM de
Intel Costa Rica y dependerán de las condiciones de acceso, autorización y
disponibilidad durante el periodo de ejecución. Si circunstancias fuera del
control del proyecto impiden realizarlas en TEM, se utilizará un laboratorio de
la Escuela de Ingeniería Electrónica del Tecnológico de Costa Rica, Campus
Tecnológico Local San Carlos.
\item El prototipo desarrollado tendrá fines académicos y 
no pretende sustituir sistemas comerciales certificados para el control ESD.
\item El diseño de un circuito propio para la medición de resistencia
cuerpo-tierra no forma parte del alcance del proyecto.
\item No se solicitará acceso a bases de datos empresariales con las
identificaciones de los usuarios del laboratorio TEM de Intel Costa Rica. Por
lo tanto, cada participante deberá registrar su nombre e identificación la
primera vez que utilice el prototipo.
\item El prototipo no medirá la cantidad de carga electrostática, el campo
electrostático ni la tensión eléctrica almacenada en el usuario.
\end{itemize}

\subsection{Consideraciones de privacidad}

Debido a que el prototipo contempla reconocimiento facial como mecanismo
complementario de identificación, se consideran medidas para proteger la
privacidad de los usuarios durante la validación del sistema. El reconocimiento
facial no será un requisito obligatorio para utilizar el prototipo, ya que la
lectura de código de barras se establece como el mecanismo principal de
identificación.

El registro facial de cada usuario se realizará únicamente con autorización
previa. En caso de que un usuario no autorice el uso de reconocimiento facial,
el sistema podrá operar mediante la lectura de código de barras. De esta forma,
el reconocimiento facial se mantiene como una función complementaria y no como
una condición obligatoria para ejecutar la prueba ESD.

La información asociada a los usuarios, como nombres, identificaciones
empresariales o datos relacionados con el registro facial, será utilizada
únicamente para validar el funcionamiento del prototipo. Estos datos no serán
publicados en el documento final ni compartidos fuera del entorno autorizado de
validación.

Los reportes históricos generados localmente por el prototipo podrán incluir el
nombre y el identificador de los usuarios autorizados. No obstante, cualquier
resultado incorporado al informe académico final se presentará de forma agregada
o anonimizada.

Cuando sea necesario presentar resultados en el informe final, la información se
mostrará de forma agregada o anonimizada. Por ejemplo, se podrán reportar
cantidades de ciclos ejecutados, tasas de lectura correcta, tasas de
reconocimiento y cantidad de registros completos, sin incluir nombres,
fotografías ni identificaciones empresariales de los participantes.

Asimismo, la información utilizada por el prototipo será almacenada localmente y
no se contempla su envío a servicios en la nube, bases de datos corporativas o
sistemas externos. Esta decisión reduce la exposición de datos personales y
mantiene el alcance del proyecto dentro de una validación académica controlada.

Finalmente, el uso del reconocimiento facial se limitará a la validación técnica
del subsistema de identificación. La finalidad de esta función no es realizar
vigilancia, control de asistencia ni seguimiento permanente de usuarios, sino
evaluar si puede utilizarse como mecanismo complementario dentro del flujo de
verificación ESD propuesto.
